\documentclass[12pt]{article}
\usepackage[utf8]{inputenc}
\usepackage[T1]{fontenc}
\usepackage[margin=1in]{geometry}
\usepackage{setspace}
\usepackage{xcolor}
\usepackage[colorlinks=true, urlcolor=blue!60!black, linkcolor=blue!60!black]{hyperref}
\setlength{\parindent}{0em}
\setlength{\parskip}{0.6em}
\renewcommand{\baselinestretch}{1.0}
\begin{document}
\begin{center}
{\LARGE \textit{Wisconsin v. Mitchell}, 508 U.S. 476 (1993)\par}
\vspace{0.5em}
{\large Case Brief --- Two-Page Edition\par}
\end{center}

\section{Memory Jogger}
A hate-crime penalty enhancement for selecting a victim by race punishes conduct, not speech, and is consistent with the First Amendment (\href{https://www.courtlistener.com/opinion/112879/wisconsin-v-mitchell/}{CourtListener}).

\section{Detailed Case Facts}
After discussing a scene from \textit{Mississippi Burning}, Todd Mitchell asked his group, ``Do you all feel hyped up to move on some white people?'' and sent them after a passing white boy, who was beaten into a four-day coma. Wisconsin's \S~939.645 raised the maximum penalty when a defendant selects a victim because of race or another listed trait.

\section{Procedural History}
A jury convicted Mitchell of aggravated battery and found racial selection, raising his maximum from two to seven years; he received four. The Wisconsin Supreme Court struck down the enhancement as punishing ``offensive thought.''

\section{Judicial Votes}
Reversed 9--0. Chief Justice Rehnquist wrote for a unanimous Court; no separate opinions.

\section{Holding}
A State may enhance the penalty for unprotected criminal conduct when the victim is selected because of a protected trait, and a defendant's prior speech may be used to prove motive or intent.

\section{Analysis of Opinions}
\textbf{Majority:} Physical assault is not expressive conduct, and sentencing has long weighed motive; \textit{Dawson v. Delaware} bars only punishing ``abstract beliefs'' unrelated to the crime. Like Title VII, the statute targets discriminatory conduct. Unlike the speech ordinance in \textit{R.A.V.}, it reaches conduct, and bias crimes inflict greater harm, including retaliation and community unrest. The claimed chilling effect was too speculative.

\section{Examples: Future Applications}
\textbf{Same-Side:} enhancing an arson sentence for burning a church because of its members' religion. \textbf{Opposite-Side:} jailing people for shouting epithets with no underlying crime (\textit{R.A.V.}). \textbf{Fence-Sitter:} a protester who screams slurs and shoves a sign.

\section{Critique}
Civil libertarians such as Susan Gellman argue the extra penalty punishes only motive and lets legislatures evade \textit{R.A.V.} Civil-rights scholars such as Frederick Lawrence answer that bias violence terrorizes whole communities.

\section{Key Quotations}
\begin{quote}
``The First Amendment, moreover, does not prohibit the evidentiary use of speech to establish the elements of a crime or to prove motive or intent.''
\end{quote}
\noindent --- Chief Justice Rehnquist, Opinion of the Court, 508 U.S. at 489.

\end{document}
